\section{Limitations}

This study is a focused pilot rather than a general treatment of recursive training. The primary evidence covers one licensed text domain and one small screening-scale language model. Behaviour at this scale may not transfer to frontier pretraining, other model families, multimodal systems, instruction tuning, or production data pipelines. A consistent result across the executed chains supports only the frozen domain, model, horizon, decoding rule, and budget regime, and the result reported here is consistent across them.

The monitored modes are constructed rather than naturally complete. Semantic clusters or available labels may merge distinct subpopulations, split a coherent mode, or omit a socially or scientifically important tail. A targeted policy can therefore be efficient relative to its monitor while being harmful relative to the full human distribution. The monitoring-omission intervention probes one form of this failure, but it cannot enumerate all misspecification. Mode-specific statistics are also noisy at small sample sizes, and an adaptive policy may spend budget in response to estimation error rather than genuine degradation.

Exact budget matching creates its own simplifications. Counting optimizer-consumed tokens is more auditable than counting documents, but equal token counts do not guarantee equal informational value, annotation cost, licensing cost, duplication, or compute. Repeated human examples may be less valuable than fresh examples even when their token presentations are counted equally. Conversely, synthetic examples can vary widely in quality and generation cost. The experiment does not establish which human-data policy is best under any economic accounting; it compares four allocation rules at equal token cost, which is a narrower question.

Compute limits constrain both the number of independent chains and the breadth of baselines, and they bind unevenly. Generations within a chain are dependent observations and cannot substitute for additional seeds. Between-chain variance turned out small enough that the frozen five-seed set exceeds what the preregistered threshold requires, so chain count is not what limits this study; the breadth of implemented baselines is. Screening a schedule or selector on the same small system can also overfit policy hyperparameters. Failed and invalid runs are retained rather than discarded, objective invalidation rules are applied, and no divergent chain was replaced.

The literature contains strong alternatives: accumulation~\cite{gerstgrasser2024inevitable}, fixed real-data mixing, correction or verification~\cite{gillman2024selfcorrecting}, detector-based resampling, and dynamic domain-mixture optimization~\cite{xie2023doremi}. Three of the seven comparators this work predeclared were not implemented, so any of these may offer an advantage this experiment could not have detected, and every omitted strong baseline narrows what the reported comparison covers. The joint policy in fact offered no detectable advantage over the strongest baseline that \emph{was} implemented, which bounds the claim from the other side as well. In particular, this work does not claim that synthetic data inevitably causes collapse, that human data is always superior, that time-and-mode allocation is new in general, or that any policy measured here is the best available.

Six limitations attach specifically to the executed grid and belong here rather than in the results. First, the primary contrast is a null and not a demonstration of inferiority: the interval lies inside the practical equivalence region, which licenses the statement that joint allocation does not improve on the strongest baseline by a practically meaningful margin, and licenses no statement that it is worse. Second, three of the seven predeclared comparators were never implemented --- accumulation or fixed-fraction mixing, detector-based or importance-resampling selection, and oracle mode information as a non-deployable upper bound --- so ``strongest eligible baseline'' means strongest among four. The absent oracle is the consequential one: without it, an equivalence between two policies carries no information about how much of the achievable gain either captures, and a reader should not infer that selection-only sits near a ceiling. Third, the mechanism by which the earlier additive assembly rule made totals diverge across arms is still not established --- that run recorded chain totals but not per-generation token accounting, so whether it arose from block packing or from divergent generated text is unresolved. Displacement removes the divergence by construction and is validated to that extent, but it is a structural remedy adopted without the diagnosis. Fourth, all twenty-five chains classify \texttt{valid\_with\_limitation} rather than \texttt{valid}, on two standing grounds: the realised token ledger was compared against the frozen budget rather than recomputed from batch records, and near-duplicate coverage was not rescanned at validation time. Fifth, the retained artifacts record how much human data each policy consumed but not \emph{when} within the chain or on which monitored modes: the per-generation allocation records live only in the chain checkpoints, which were not archived before the compute was released. The joint and selection-only arms consumed identical lifetime totals, but whether the joint policy reached that total by a different temporal or distributional route cannot be determined from what was kept. This matters for reading the primary null: an equivalence produced by two policies that allocated differently is a statement about allocation strategy, and one produced by two policies that converged on the same allocation is a statement about the joint rule. We cannot distinguish them here. Sixth, the practical effect threshold that every interval in this paper is read against was frozen before outcomes existed and has been checked against measurement, but has never had the independent statistics review the project's own decision log calls for; the novelty assessment is likewise internal, with a literature search that carries a date horizon. Because the primary finding is an equivalence, it rests on that threshold more directly than a positive finding would.

The empirical claims are confined to one operating point. A published positive control has been reproduced and is reported as such; the grid executed with both fairness axes holding and produced certifiable artifacts; and the comparisons among the treatment families are valid for the frozen domain, model, horizon, decoding rule, and budget regime, and for no other. The method, exact domain, monitoring map, effect threshold, and powered seed count remain subject to preregistration gates. Conclusions report null, harmful, or ambiguous outcomes where that is what the immutable artifacts support, and the primary outcome reported here is a null.
